\documentclass[aspectratio=169]{beamer}
\usetheme{Boadilla}
\usecolortheme{seahorse}
\usepackage[utf8]{inputenc}
\usepackage[spanish]{babel}
\usepackage{graphicx}
\usepackage{amsmath}

\title[Introducción y Morfología]{Clase 1: Introducción a la Robótica y Morfología de Manipuladores}
\subtitle{Robótica (IMT-342)}
\author[B. Quiroga Turdera]{M.Sc. Ing. Bernardo Quiroga Turdera}
\institute[UCB Tarija]{Universidad Católica Boliviana ``San Pablo'' \\ Sede Tarija}
\date{4 de agosto de 2026}

\begin{document}

\begin{frame}
    \titlepage
\end{frame}

\begin{frame}{Encuadre de la Asignatura}
    \textbf{Competencia General:} \\
    Diseñar, simular y/o implementar sistemas robóticos para reemplazar o colaborar a la acción humana en tareas complejas, repetitivas y de alto riesgo. %[cite: 1]
    
    \vspace{0.5cm}
    \textbf{Sistema de Evaluación:} %[cite: 1]
    \begin{itemize}
        \item \textbf{50\% Evaluación Continua:} Hitos del Proyecto Final Integrador (PFI), laboratorios y exámenes cortos.
        \item \textbf{50\% Evaluación Final:} Defensa de la celda robótica operando de manera 100\% autónoma.
    \end{itemize}
    \vspace{0.3cm}
    \textit{Nota: La habilitación para la evaluación final requiere un mínimo de 50/100 puntos en la evaluación continua.} %[cite: 1]
\end{frame}

\begin{frame}{El Reto del Semestre (PFI)}
    \begin{block}{Celda de Manufactura Flexible Inteligente}
        Diseño, simulación e implementación de una celda de manufactura con el \textbf{ABB IRB 120} guiado por visión artificial.
    \end{block}
    
    \textbf{Restricciones de Ingeniería:}
    \begin{itemize}
        \item Precisión exigida: Error absoluto lineal $< 1.5$ mm.
        \item Tiempo de ciclo: $< 8$ segundos por pieza.
        \item \textbf{Regla Anti-Caja Negra:} Prohibido usar librerías iterativas (MoveIt!) para cinemática. Todos los algoritmos matemáticos deben codificarse desde cero en Python.
    \end{itemize}
\end{frame}

\begin{frame}{Morfología de un Manipulador Industrial}
    \textbf{Cadena Cinemática Abierta:}
    Secuencia de eslabones rígidos unidos por articulaciones rotacionales (Juntas de Revolución $R$).
    
    \vspace{0.3cm}
    \textbf{Anatomía del ABB IRB 120 (6 GDL):} %[cite: 1]
    \begin{itemize}
        \item \textbf{Base y Hombro (Juntas 1 y 2):} Posicionamiento espacial grueso.
        \item \textbf{Codo (Junta 3):} Alcance en el plano vertical.
        \item \textbf{Muñeca Esférica (Juntas 4, 5 y 6):} Orientación del efector final. Los tres ejes intersectan en un único punto común ($p_w$), clave para el Desacoplamiento de Pieper.
    \end{itemize}
\end{frame}

\begin{frame}{Actuación y Transmisión}
    \begin{columns}
        \column{0.5\textwidth}
        \textbf{Servomotores AC Brushless:}
        \begin{itemize}
            \item Control de corriente preciso.
            \item Codificadores absolutos en cada junta.
        \end{itemize}
        
        \vspace{0.3cm}
        \textbf{Reductores Armónicos (Harmonic Drives):}
        \begin{itemize}
            \item Cero juego mecánico (backlash).
            \item Alta relación de reducción en volumen compacto.
        \end{itemize}
        
        \column{0.5\textwidth}
        \textbf{Carga Útil y Espacio de Trabajo:}
        \begin{itemize}
            \item Alcance máximo: $580$ mm.
            \item Carga útil nominal (Payload): $3$ kg.
            \item Repetibilidad de posición: $0.01$ mm.
        \end{itemize}
    \end{columns}
\end{frame}

\begin{frame}{Próximos Pasos}
    \begin{block}{Para la próxima semana}
        Dado que este jueves es feriado nacional (6 de agosto), nos reencontramos el próximo martes. 
        \vspace{0.2cm}
        
        \textbf{Requisito obligatorio:}
        Traer laptops con Ubuntu 24.04 LTS y ROS instalados para comenzar a trabajar con los archivos CAD y la geometría del Hito 1.
    \end{block}
\end{frame}

\end{document}